% 核心结论：负载相位改变电压调整
% 图原型：纯几何／对比。物理网络按本题节点、端口和符号构建。
\documentclass[tikz,border=7pt]{standalone}
\usepackage{fontspec,xeCJK,amsmath,amssymb}
\IfFontExistsTF{Times New Roman}{\setmainfont{Times New Roman}}{\setmainfont{TeX Gyre Termes}}
\IfFontExistsTF{Heiti SC}{\setCJKmainfont[AutoFakeBold=1.5]{Heiti SC}}{\setCJKmainfont[AutoFakeBold=1.5]{FandolHei-Regular}}
\usetikzlibrary{arrows.meta,calc,angles,quotes}
\definecolor{ink}{HTML}{243745}
\definecolor{blue}{HTML}{275DAB}
\definecolor{green}{HTML}{177C63}
\definecolor{red}{HTML}{BD4935}
\definecolor{purple}{HTML}{715393}
\begin{document}
\begin{tikzpicture}[font=\fontsize{9.5}{12}\selectfont,text=ink,
 v/.style={line width=1.1pt,-{Stealth[length=2.4mm,width=1.5mm]}},
 axis/.style={ink!65,line width=.45pt,-{Stealth[length=1.7mm,width=1.1mm]}},
 aux/.style={ink!40,densely dashed,line width=.5pt},
 lab/.style={fill=white,inner sep=2pt},
 card/.style={draw=blue!25,fill=blue!3,rounded corners=3pt,text width=6cm,align=left,inner sep=10pt}]
\path[use as bounding box] (-8,-3.35) rectangle (8,3.65);
\node[font=\fontsize{10.5}{13}\selectfont\bfseries,text width=15cm,align=center] at (0.00000,3.22000) {M6-061　负载相位改变电压调整};
\begin{scope}[shift={(-4,0)}]
\draw[ink,thick] (-2.60000,1.50000) -- (-2.60000,0.40000);
\draw[ink,thick] (-2.60000,-0.40000) -- (-2.60000,-1.50000);
\draw[ink,thick] (-2.6,0) circle (.4);
\node[] at (-2.60000,0.00000) {$\sim$};
\node[] at (-3.22000,0.00000) {$U_1$};
\draw[ink,thick] (-2.60000,1.50000) -- (-1.40000,1.50000);
\draw[ink,thick] (-2.60000,-1.50000) -- (2.80000,-1.50000);
\draw[ink,thick] (-1.40000,1.50000) -- (-1.40000,0.42000);
\draw[ink,thick] (-1.40000,-0.42000) -- (-1.40000,-1.50000);
\begin{scope}[shift={(-1.40000,0.00000)},rotate=-90.0]
\draw[ink,thick,fill=white] (-0.42,-.14) rectangle (0.42,.14);
\end{scope}
\node[lab,anchor=west] at (-1.12000,0.00000) {$R_c$};
\draw[ink,thick] (-1.40000,1.50000) -- (-0.45000,1.50000);
\draw[ink,thick] (-0.45000,1.50000) -- (-0.45000,0.42000);
\draw[ink,thick] (-0.45000,-0.42000) -- (-0.45000,-1.50000);
\begin{scope}[shift={(-0.45000,0.00000)},rotate=-90.0]
\draw[ink,thick] plot[domain=-0.42:0.42,samples=70] (\x,{.20*abs(sin(720*(\x+0.42)/0.84))});
\end{scope}
\node[lab,anchor=west] at (-0.03000,0.00000) {$jX_m$};
\draw[ink,thick] (-0.45000,1.50000) -- (0.75500,1.50000);
\draw[ink,thick] (1.59500,1.50000) -- (2.80000,1.50000);
\begin{scope}[shift={(1.17500,1.50000)},rotate=0.0]
\draw[ink,thick,fill=white] (-0.42,-.14) rectangle (0.42,.14);
\end{scope}
\node[lab] at (1.17500,1.92000) {$Z_{eq}$};
\draw[ink,thick] (2.80000,1.50000) -- (2.80000,0.42000);
\draw[ink,thick] (2.80000,-0.42000) -- (2.80000,-1.50000);
\begin{scope}[shift={(2.80000,0.00000)},rotate=-90.0]
\draw[ink,thick,fill=white] (-0.42,-.14) rectangle (0.42,.14);
\end{scope}
\node[lab,anchor=east] at (2.38000,0.00000) {$Z_L'$};
\node[] at (0.00000,-2.40000) {折算至初级：全部电压、电流与阻抗取同一侧。};
\end{scope}
\begin{scope}[shift={(4,0)}]
\node[card] at (0.00000,1.15000) {$R_{eq}=1.2\Omega,\ X_{eq}=1.6\Omega$};
\end{scope}
\begin{scope}[shift={(4,0)}]
\node[card] at (0.00000,-1.15000) {$\Delta U\simeq I(R\cos\phi\pm X\sin\phi)$};
\end{scope}
\node[text width=15cm,align=center] at (0.00000,-3.00000) {折算负载端电压作为实轴，滞后和超前电流使漏抗压降的实轴投影取不同符号。};

\end{tikzpicture}\end{document}